\scanpage{15}
En vertu de la déf. (12), le deuxième mb est $f\alpha'(f(x)\wedge f(y))$, de sorte qu'il suffit de prouver
\[
 [x,y]_{\varphi'}=\deleted{\unsure{\alpha'f(f(x)}}\ \alpha'(f(x)\wedge f(y)) \tag*{(50)}
\]
ce qui n'est autre que la définition (40 bis).

\uline{NB} En général, $f':E'\to E$ n'est pas compatible aux structures d'Algèbres, $=$ si $f$ est un isomorphisme i.e. $\Delta$ une unité. Dans ce cas, on a $f'=\Delta f^{-1}$ (45), \deleted{\unsure{donc}} \unsure{or|} $f^{-1}$ est \deleted{\unsure{un}} isom. d'Algèbres comme inverse d'un isomorphisme, i.e. on a
\[
 f^{-1}([x',y'])=[f^{-1}(x'),f^{-1}(y')]
\]
Si $f'=\Delta f^{-1}$ \unsure{était} un hom d'algèbre, la \deleted{\unsure{même}} formule précédente \unsure{resterait} valable en y remplaçant $f'$ par $\Delta f^{-1}$, i.e. on aurait à la fois $a=b$ et $\Delta a=\Delta^{\unsure{\prime|2}} b$ i.e. $a=\Delta b$, donc $\Delta a=a$. \deleted{Si les} On en conclut aisément (comme les $[x',y']$ engendrent $E'$ localement ...) que $\Delta=1$, \deleted{Donc} i.e. $f'=f^{-1}$. Donc dès que $\Delta\ne1$, $\Delta$ inversible, $f'$ n'est \uline{\unsure{p=|pas}} un hom d'Algèbres....

% archive sheet number: 14
